\documentclass[11pt,a4paper]{article}

\usepackage[utf8]{inputenc}
\usepackage[T1]{fontenc}
\usepackage{amsmath,amssymb,amsfonts,amsthm}
\usepackage{mathtools}
\usepackage{geometry}
\usepackage{booktabs}
\usepackage{hyperref}
\usepackage{enumitem}

\geometry{margin=1in}
\hypersetup{
    colorlinks=true,
    linkcolor=blue,
    citecolor=blue,
    urlcolor=blue,
    pdftitle={The Softmax Inundation Bound},
    pdfauthor={Shaurya Sanyal}
}

\theoremstyle{plain}
\newtheorem{theorem}{Theorem}
\newtheorem{lemma}{Lemma}
\newtheorem{corollary}{Corollary}
\newtheorem{proposition}{Proposition}

\theoremstyle{definition}
\newtheorem{definition}{Definition}
\newtheorem{assumption}{Assumption}

\theoremstyle{remark}
\newtheorem{remark}{Remark}

\title{%
  \textbf{The Softmax Inundation Bound}\\[0.4em]
  \large Attention Dilution in Long-Context Transformers and\\
  Generational Context as a Volume-Bounding Remedy
}

\author{%
  \textbf{Shaurya Sanyal}\\
  \textit{Team Auralis}\\
  \texttt{shauryasanyal@gmail.com}
}

\date{October 2026}

\begin{document}

\maketitle

% ---------------------------------------------------------------------------
\begin{abstract}
\noindent
We isolate and formalize a simple but under-appreciated consequence of softmax
normalization in transformer self-attention: when a fixed set of \emph{invariant}
tokens competes with a growing set of distractor tokens, the aggregate attention
mass on the invariant set is squeezed by the normalization denominator. We call
this \emph{softmax inundation}. Under the standard bounded-projection-energy
assumption we prove
\[
  A_{i,\mathrm{inv}} \;\le\; \frac{\kappa\, L_u}{L_u + |\mathcal{S}_{\mathrm{noise}}|}
  \;=\; \mathcal{O}\!\left(\frac{L_u}{|\mathcal{S}_{\mathrm{noise}}|}\right),
  \qquad \kappa = e^{2M_QM_K/\sqrt{d_k}},
\]
so that the invariant share decays at least as fast as the inverse of the
distractor count.

Because the inequality is conditional on uniform score bounds and can be
uninformative when score contrast is large, we do not present it as a new law of
attention. Its value is diagnostic: it isolates \emph{token volume} as the
variable a system designer can actually control.

We then introduce \textbf{Generational Context}, a three-tier agent memory
architecture (Gen-0 ephemeral tool output with TTL distillation, Gen-1 working
memory paged by Attention-Weighted Eviction, Gen-2 permanently pinned
invariants), and prove that under a strict bound $C_{\mathrm{bound}}$ on active
non-invariant volume, invariant attention admits a \emph{turn-independent}
lower bound,
\[
  A_{i,\mathrm{inv}}^{(\mathrm{GC})} \;\ge\;
  \frac{L_u}{L_u + \kappa\, C_{\mathrm{bound}}} \;>\; 0
  \qquad \text{for all } t \ge 1 .
\]
This is the substantive claim: bounding volume removes horizon-dependence from
the worst case. It is a guarantee about attention allocation, not about task
success. We report algorithmic-simulation results, describe a reference
implementation with cryptographic intent auditing, and state precisely what
remains unverified.
\end{abstract}

% ---------------------------------------------------------------------------
\section{Introduction}
\label{sec:intro}

Autonomous agents built on transformer language models typically treat the
context window as an append-only transcript. Every tool call, every stack trace,
every intermediate hypothesis is appended and never removed. This is convenient
to implement and expensive to sustain.

Two things are routinely conflated in discussion of this practice:

\begin{enumerate}[label=(\roman*)]
  \item \textbf{Capacity.} Whether the transcript still fits in the window.
  \item \textbf{Allocation.} Whether attention still lands on the tokens that
        matter.
\end{enumerate}

Work on long-context models has largely addressed (i): windows grow, and
retrieval quality within a window is measured by benchmarks such as RULER
\cite{hsieh2024ruler} and by position-sensitive studies such as
\emph{Lost in the Middle} \cite{liu2023lost}. But (ii) is a separate phenomenon,
and it is governed by the softmax denominator rather than by the window size.

Our contributions are deliberately modest and, we hope, clearly separated:

\begin{itemize}
  \item \S\ref{sec:inundation} states and proves the \emph{Softmax Inundation
        Bound}. It is elementary; we claim only that it is a useful way to
        isolate the mechanism, not that the mechanism was previously unknown.
  \item \S\ref{sec:gc} introduces Generational Context and proves the
        \emph{Generational Retention Guarantee}, a horizon-independent lower
        bound on invariant attention under bounded active volume. This is the
        result we consider the paper's actual contribution.
  \item \S\ref{sec:awe}--\ref{sec:impl} describe the eviction policy, the
        reference implementation, and a simulation harness.
  \item \S\ref{sec:limits} is an explicit limitations section. We regard it as
        part of the contribution rather than as boilerplate.
\end{itemize}

\paragraph{What this paper does not claim.}
It does not claim that attention dilution causes reasoning failure in any
particular model; residual streams, induction heads, and tool use can all
compensate. It does not claim empirical gains on live models; we have not run
those experiments. It does not claim novelty for softmax competition itself.

% ---------------------------------------------------------------------------
\section{Preliminaries}
\label{sec:prelim}

\subsection{Interaction transcript}

Let an agent interact with an environment over discrete turns $t \in
\mathbb{N}$. At turn $t$ the context is
\begin{equation}
  \mathcal{H}_t
  = \bigl(\mathbf{u},\;
    \mathbf{a}_1,\mathbf{o}_1,\mathbf{r}_1,\;
    \dots,\;
    \mathbf{a}_{t-1},\mathbf{o}_{t-1},\mathbf{r}_{t-1}\bigr),
\end{equation}
where $\mathbf{u} \in \mathcal{V}^{L_u}$ holds the user objective and system
invariants, $\mathbf{a}_k$ are actions, $\mathbf{o}_k$ are raw environmental
observations, and $\mathbf{r}_k$ are reasoning traces. The context length is
\begin{equation}
  N_t = L_u + \sum_{k<t}\bigl(L_a^{(k)} + L_o^{(k)} + L_r^{(k)}\bigr).
\end{equation}

\subsection{Scaled dot-product attention}

For query position $i$ and key position $j$,
\begin{equation}
  A_{i,j}
  = \frac{\exp\!\bigl(s_{i,j}\bigr)}{\sum_{m=1}^{N_t}\exp\!\bigl(s_{i,m}\bigr)},
  \qquad
  s_{i,j} = \frac{\mathbf{q}_i^{\top}\mathbf{k}_j}{\sqrt{d_k}},
  \label{eq:attn}
\end{equation}
with $\mathbf{q}_i = \mathbf{X}_i W_Q$ and $\mathbf{k}_j = \mathbf{X}_j W_K$.

\subsection{Partition of context}

Fix a turn $t$ and partition the index set $J = \{1,\dots,N_t\}$ into three
mutually disjoint subsets:
\begin{align}
  \mathcal{S}_{\mathrm{inv}} &: \text{invariant tokens}, &
  |\mathcal{S}_{\mathrm{inv}}| &= L_u, \\
  \mathcal{S}_{\mathrm{noise}} &: \text{raw observations}, &
  |\mathcal{S}_{\mathrm{noise}}| &= \textstyle\sum_k L_o^{(k)}, \\
  \mathcal{S}_{\mathrm{work}} &: \text{actions and reasoning}, &
  |\mathcal{S}_{\mathrm{work}}| &= \textstyle\sum_k \bigl(L_a^{(k)}+L_r^{(k)}\bigr).
\end{align}
Write $A_{i,\mathrm{inv}} = \sum_{j \in \mathcal{S}_{\mathrm{inv}}} A_{i,j}$.

\begin{assumption}[Bounded projection energy]
\label{assump:bounded}
There exist $M_Q, M_K > 0$ such that $\|\mathbf{q}_i\|_2 \le M_Q$ and
$\|\mathbf{k}_j\|_2 \le M_K$ for all $i,j$. Consequently, by Cauchy--Schwarz,
\begin{equation}
  -B \;\le\; s_{i,j} \;\le\; B,
  \qquad B = \frac{M_Q M_K}{\sqrt{d_k}} .
\end{equation}
Define the dynamic range factor $\kappa = e^{2B}$.
\end{assumption}

\begin{remark}
Assumption~\ref{assump:bounded} is weak but not innocuous. Weight matrices in
trained transformers are finite, so it always holds for \emph{some} $M_Q, M_K$;
the question is whether the resulting $\kappa$ is small enough for the bound to
say anything. See \S\ref{sec:limits}.
\end{remark}

% ---------------------------------------------------------------------------
\section{The Softmax Inundation Bound}
\label{sec:inundation}

\begin{theorem}[Softmax Inundation Bound]
\label{thm:inundation}
Under Assumption~\ref{assump:bounded}, for every query position $i$ and every
turn $t$,
\begin{equation}
  A_{i,\mathrm{inv}}
  \;\le\;
  \frac{\kappa\, L_u}{L_u + |\mathcal{S}_{\mathrm{noise}}|},
  \qquad \kappa = e^{2M_QM_K/\sqrt{d_k}} .
  \label{eq:inundation}
\end{equation}
In particular, holding $L_u$ and the scoring regime fixed,
\begin{equation}
  A_{i,\mathrm{inv}} = \mathcal{O}\!\left(\frac{L_u}{|\mathcal{S}_{\mathrm{noise}}|}\right)
  \quad \text{as } |\mathcal{S}_{\mathrm{noise}}| \to \infty .
\end{equation}
\end{theorem}

\begin{proof}
Bounds on $s_{i,j}$ give $e^{-B} \le e^{s_{i,j}} \le e^{B}$ for all $j$. Hence
\begin{equation}
  \sum_{j \in \mathcal{S}_{\mathrm{inv}}} e^{s_{i,j}} \;\le\; L_u e^{B},
\end{equation}
and, dropping the non-negative $\mathcal{S}_{\mathrm{work}}$ terms from the
denominator of \eqref{eq:attn},
\begin{equation}
  \sum_{m=1}^{N_t} e^{s_{i,m}}
  \;\ge\; \bigl(L_u + |\mathcal{S}_{\mathrm{noise}}|\bigr) e^{-B}.
\end{equation}
Dividing,
\begin{equation}
  A_{i,\mathrm{inv}}
  \;\le\;
  \frac{L_u e^{B}}{\bigl(L_u + |\mathcal{S}_{\mathrm{noise}}|\bigr) e^{-B}}
  \;=\;
  \frac{e^{2B} L_u}{L_u + |\mathcal{S}_{\mathrm{noise}}|}.
\end{equation}
With $\kappa = e^{2B}$ this is \eqref{eq:inundation}.
\end{proof}

\begin{remark}[Reading the bound correctly]
\label{rem:reading}
Equation~\eqref{eq:inundation} is an \emph{upper} bound on the invariant share.
It is informative when $\kappa$ is moderate; it becomes vacuous as soon as
$\kappa L_u \gtrsim L_u + |\mathcal{S}_{\mathrm{noise}}|$, at which point the
right-hand side exceeds $1$ and the inequality carries no information. This is
not a defect of the proof but a genuine limitation of the argument: bounding
norms alone does not force dilution. If distractor scores are systematically
much smaller than invariant scores, $A_{i,\mathrm{inv}} \to 1$ regardless of
$|\mathcal{S}_{\mathrm{noise}}|$.
\end{remark}

\begin{remark}[Relationship to known results]
Softmax competition of this form is implicit in any analysis of attention with a
growing key set, and explicit in streaming and long-context work
\cite{xiao2023streamingllm,liu2023lost}. We do not claim the mechanism as new.
The useful part, in our view, is the separation of \emph{volume} (which a system
controls) from \emph{score contrast} (which it largely does not), because only
the former is addressable by memory management.
\end{remark}

% ---------------------------------------------------------------------------
\section{Generational Context}
\label{sec:gc}

If Theorem~\ref{thm:inundation} identifies volume as the controllable variable,
the natural response is to control it. Generational Context borrows the shape of
generational garbage collection and of virtual-memory paging, and applies both
at the granularity of agent context rather than of KV blocks.

\begin{definition}[Generational Context]
\label{def:gc}
Active context is partitioned into three tiers:
\begin{enumerate}[label=\textbf{Gen-\arabic*}, leftmargin=*]
  \item \emph{Tenured.} Founding objectives and system invariants. Pinned;
        never evicted. Occupies $L_u$ tokens.
  \item \emph{Evaporative.} Working memory: hypotheses, plans, uncommitted
        deltas. Paged by Attention-Weighted Eviction (\S\ref{sec:awe}) under a
        strict capacity $C_1$.
  \item \emph{Ephemeral.} Raw tool output and logs. Passed through a
        distillation map $\sigma_{\mathrm{distill}}$ producing at most $K_0$
        tokens, with the raw form evicted under a turn-based TTL.
\end{enumerate}
Let $C_{\mathrm{overhead}}$ bound the structural tokens (role tags, delimiters,
section headers) required to render the window, and set
\begin{equation}
  C_{\mathrm{bound}} = C_1 + K_0 + C_{\mathrm{overhead}} .
\end{equation}
\end{definition}

\begin{theorem}[Generational Retention Guarantee]
\label{thm:retention}
Under Assumption~\ref{assump:bounded}, and given a Generational Context with
$C_{\mathrm{bound}} < \infty$,
\begin{equation}
  A_{i,\mathrm{inv}}^{(\mathrm{GC})}
  \;\ge\;
  \frac{L_u}{L_u + \kappa\, C_{\mathrm{bound}}}
  \;>\; 0
  \qquad \text{for all } t \ge 1 .
  \label{eq:retention}
\end{equation}
In particular the bound is independent of the interaction horizon $t$.
\end{theorem}

\begin{proof}
By Definition~\ref{def:gc}, at every turn the number of active non-invariant
tokens satisfies
\begin{equation}
  \bigl|\mathcal{S}_{\mathrm{noninv}}^{(t)}\bigr|
  \;\le\; C_1 + K_0 + C_{\mathrm{overhead}}
  \;=\; C_{\mathrm{bound}}
\end{equation}
independently of $t$. Set
\begin{equation}
  R_t
  = \frac{\sum_{m \notin \mathcal{S}_{\mathrm{inv}}} e^{s_{i,m}}}
         {\sum_{j \in \mathcal{S}_{\mathrm{inv}}} e^{s_{i,j}}},
  \qquad\text{so that}\qquad
  A_{i,\mathrm{inv}} = \frac{1}{1 + R_t}.
\end{equation}
The numerator has at most $C_{\mathrm{bound}}$ terms, each at most $e^{B}$; the
denominator has $L_u$ terms, each at least $e^{-B}$. Therefore
\begin{equation}
  R_t
  \;\le\;
  \frac{C_{\mathrm{bound}}\, e^{B}}{L_u\, e^{-B}}
  \;=\;
  \frac{\kappa\, C_{\mathrm{bound}}}{L_u}.
\end{equation}
Since $R \mapsto 1/(1+R)$ is decreasing on $[0,\infty)$,
\begin{equation}
  A_{i,\mathrm{inv}}
  \;\ge\;
  \frac{1}{1 + \kappa C_{\mathrm{bound}}/L_u}
  \;=\;
  \frac{L_u}{L_u + \kappa C_{\mathrm{bound}}} .
\end{equation}
Because $C_{\mathrm{bound}}$ is fixed by the memory controller, the right-hand
side does not depend on $t$.
\end{proof}

\begin{corollary}[No horizon decay]
\label{cor:no-decay}
Under the hypotheses of Theorem~\ref{thm:retention},
$\inf_{t \ge 1} A_{i,\mathrm{inv}}^{(t)} > 0$. That is, unbounded softmax
inundation cannot occur, in the worst case, once active non-invariant volume is
bounded.
\end{corollary}

\begin{remark}[What the guarantee is and is not]
The bound is a \emph{worst-case} statement under uniform score assumptions. It
does not say invariant attention is large; it says it cannot be driven to zero by
volume alone. Task success additionally requires correct retrieval, valid tool
use, and sound multi-step reasoning, none of which this theorem addresses. An
adversary who controls distractor content rather than distractor count can
violate the spirit of the result by inflating $s_{i,m}$ while staying inside
Assumption~\ref{assump:bounded}; see \S\ref{sec:limits}.
\end{remark}

% ---------------------------------------------------------------------------
\section{Attention-Weighted Eviction}
\label{sec:awe}

Theorem~\ref{thm:retention} is a statement about a controller that actually
maintains the volume bound. We use an adaptation of ARC \cite{megiddo2003arc},
which already solves the \emph{recency versus frequency} allocation problem for
paged memory.

The cache maintains a recency list $T_1$, a frequency/saliency list $T_2$, ghost
registries $B_1, B_2$ recording recently evicted keys, and an adaptation
parameter $p$ targeting the size of $T_1$. Victim selection within a list uses an
\emph{Information Density Score}
\begin{equation}
  \mathrm{IDS}(E)
  = \frac{\mathrm{AttnWeight}(E) \cdot 100}{\mathrm{Entropy}(E) + 10^{-4}} ,
  \label{eq:ids}
\end{equation}
with tier-based priors for entries whose attention statistics have not yet been
measured. The cache enforces the post-condition
\begin{equation}
  \sum_{e \in T_1 \cup T_2} \mathrm{tokens}(e) \;\le\; C
  \qquad \text{after every operation,}
\end{equation}
truncating any single entry larger than $C$ before insertion.

Two implementation notes matter for correctness. First, promoting an entry from
$T_1$ to $T_2$ must evict \emph{other} entries if the promotion grows the
footprint; a naive implementation can select the newly promoted entry as its own
victim. Second, ghost registries must be bounded, or the eviction bookkeeping
grows without limit even though the active cache does not.

% ---------------------------------------------------------------------------
\section{Simulation Results}
\label{sec:sim}

We report results from a deterministic simulation harness (\texttt{ccb\_bench.py})
that models three arms over scripted multi-turn scenarios:
flat append-only, periodic compaction, and Generational Context. Attention
quantities are computed with 128-dimensional hash-projected embeddings and an
explicit softmax, not with trained transformer weights. Table~\ref{tab:ccb}
reports final-turn means across three scenarios.

\begin{table}[ht]
\centering
\small
\begin{tabular}{lcccc}
\toprule
\textbf{Arm} & \textbf{Active tokens} & \textbf{CSR} & \textbf{Attn(inv)} & \textbf{Cost (sim.)} \\
\midrule
Flat append-only      & 5{,}099 & 0.0059 & 0.0425 & \$0.0459 \\
Periodic compaction   & 3{,}340 & 0.0077 & 0.0553 & \$0.0301 \\
\textbf{Generational Context} & \textbf{129} & \textbf{0.2000} & \textbf{0.0757} & \textbf{\$0.0012} \\
\bottomrule
\end{tabular}
\caption{CCB-001 simulation, final-turn means over three scripted scenarios.
``CSR'' is the context signal ratio; ``Attn(inv)'' is the simulated attention
share on invariant tokens. These are algorithmic proxies, not neural
measurements. Absolute values are not comparable to published transformer
results; only the ordering across arms is meaningful.}
\label{tab:ccb}
\end{table}

\paragraph{What the simulation supports.} The harness shows that the memory
policy realizes its volume target under scripted load, and that the resulting
active footprint is roughly two orders of magnitude smaller than flat
append-only. That is a statement about the controller.

\paragraph{What it does not support.} It does not demonstrate that a trained
transformer allocates attention to invariants in the same proportions, nor that
task success improves. Those require live-model experiments, which we have not
performed.

% ---------------------------------------------------------------------------
\section{Reference Implementation and Intent Auditing}
\label{sec:impl}

The architecture is implemented in the \texttt{generational\_context} package
(Python standard library only). The controller exposes tool observation,
reasoning recording, and prompt compilation; it enforces the capacity invariant
of \S\ref{sec:awe} and the TTL behaviour of Gen-0.

A separate benchmark exercises long-horizon intent conservation. In its current
form it uses Ed25519-signed intent tokens over a canonical digest of the
declared goal, values, constraints, and priorities, together with an append-only
hash-chained audit log. A proposal is accepted only if its token verifies
against the current intent digest and its signatory holds the required role.
Earlier revisions of this component used an unsalted digest of static strings,
which provided no security; that construction has been removed.

We are explicit that this is an engineering component, not a validated security
boundary. Its threat model is undocumented and its scenarios are authored by the
same team that wrote the checker.

% ---------------------------------------------------------------------------
\section{Limitations}
\label{sec:limits}

We prefer to state these plainly rather than to bury them.

\begin{enumerate}
  \item \textbf{The inundation bound is elementary and conditional.}
        Theorem~\ref{thm:inundation} follows from Cauchy--Schwarz and the
        definition of softmax. It is informative only when $\kappa$ is moderate;
        for realistic projection norms with $B$ even of order $10$, $\kappa$
        exceeds $10^{8}$ and the bound is vacuous at practical context lengths.
        The theorem is best read as a formalization of a mechanism, not as an
        explanation of observed degradation at a particular scale.

  \item \textbf{Uniform score assumptions are strong.} Both theorems assume the
        same worst-case score envelope for invariant and distractor tokens.
        Real attention scores are highly structured, and the ordering of
        $\mathcal{S}_{\mathrm{inv}}$ versus $\mathcal{S}_{\mathrm{noise}}$ in
        score space is exactly what determines the outcome. Our bounds are
        agnostic to that ordering, which is both their strength and their
        weakness.

  \item \textbf{Multi-head and layer structure is ignored.} We analyze a single
        head in isolation. Specialization across heads (induction, copy,
        positional sinks) can preserve information that a per-head bound
        suggests is lost.

  \item \textbf{No live-model validation.} All quantitative results come from a
        deterministic simulation with hash-based embeddings. We make no claims
        about behaviour of trained models, and we would regard any such claim
        drawn from these numbers as unsupported.

  \item \textbf{The retention guarantee is about allocation, not success.}
        Keeping invariants in the attention distribution is necessary at best.
        An agent can attend to its instructions and still fail.

  \item \textbf{Adversarial distractor content is out of scope.} A bounded
        number of very high-scoring distractors can dominate attention. Volume
        control does not defend against content attacks.
\end{enumerate}

\paragraph{Open empirical questions.}
Does bounded-volume management measurably reduce goal drift in long agent
sessions? Does distillation of Gen-0 preserve the error signature that made a
failure diagnosable? Does AWE's IDS criterion outperform plain recency at equal
capacity? Each is answerable with a controlled experiment and none is answered
here.

% ---------------------------------------------------------------------------
\section{Related Work}
\label{sec:related}

\paragraph{Long-context behaviour.} Position-sensitive degradation is documented
in \emph{Lost in the Middle} \cite{liu2023lost}; effective-context benchmarks
show that nominal window size overstates usable context \cite{hsieh2024ruler}.
Streaming inference with attention sinks \cite{xiao2023streamingllm} addresses
the degenerate-statistics problem of unbounded KV growth. Our concern is
different in emphasis: not how to keep generating coherently from a long stream,
but how to keep a bounded working set aligned with declared invariants.

\paragraph{Agent memory.} MemGPT \cite{packer2023memgpt} introduces virtual
context management through function calls, with paging decisions made by the
model itself. Generational Context makes those decisions with a deterministic
controller instead, trading flexibility for predictable cost. Pinning invariants
corresponds to MemGPT's always-in-context core memory; the difference is the
explicit tiering and the TTL treatment of raw observations.

\paragraph{Cache replacement.} ARC \cite{megiddo2003arc} supplies the adaptation
mechanism. Attention-score-based eviction for KV caches is well established in
the inference-efficiency literature. We apply the same idea at the granularity of
semantic context entries rather than KV blocks.

% ---------------------------------------------------------------------------
\section{Conclusion}
\label{sec:conclusion}

Softmax normalization means that adding distractor tokens competes with
invariant tokens for a fixed budget of probability mass. We state that
competition as an explicit bound, note carefully where the bound is informative
and where it collapses, and then prove the complementary result: if a controller
holds active non-invariant volume below a constant, the worst-case invariant
share is bounded away from zero independently of how long the session runs.

The second result is the one we would defend. It does not require any claim about
model internals beyond bounded projections, it identifies a quantity a system can
actually control, and it comes with an implementation whose capacity invariant is
enforced rather than assumed.

What it does not come with is evidence that any of this improves outcomes on real
tasks. That is the next experiment, and we would rather state the gap than paper
over it.

% ---------------------------------------------------------------------------
\section*{Acknowledgments}

The term \emph{Artificial Civilization Intelligence} (ACI) was coined by Raja
Dharma Tej Maddala. The architecture, proofs, implementation, and simulation
harness in this document were produced by Team Auralis.

% ---------------------------------------------------------------------------
\begin{thebibliography}{9}

\bibitem{hsieh2024ruler}
C.-P. Hsieh, S. Sun, S. Kriman, S. Acharya, D. Rekesh, F. Jia, Y. Zhang, and
B. Ginsburg.
\newblock RULER: What's the real context size of your long-context language
  models?
\newblock \emph{arXiv:2404.06654}, 2024.

\bibitem{liu2023lost}
N.~F. Liu, K. Lin, J. Hewitt, A. Paranjape, M. Bevilacqua, F. Petroni, and
P. Liang.
\newblock Lost in the middle: How language models use long contexts.
\newblock \emph{Transactions of the ACL}, 2024. arXiv:2307.03172.

\bibitem{megiddo2003arc}
N. Megiddo and D.~S. Modha.
\newblock ARC: A self-tuning, low overhead replacement cache.
\newblock In \emph{Proceedings of FAST}, 2003.

\bibitem{packer2023memgpt}
C. Packer, S. Wooders, K. Lin, V. Fang, S.~G. Patil, I. Stoica, and J.~E.
  Gonzalez.
\newblock MemGPT: Towards LLMs as operating systems.
\newblock \emph{arXiv:2310.08560}, 2023.

\bibitem{xiao2023streamingllm}
G. Xiao, Y. Tian, B. Chen, S. Han, and M. Lewis.
\newblock Efficient streaming language models with attention sinks.
\newblock In \emph{Proceedings of ICLR}, 2024. arXiv:2309.17453.

\end{thebibliography}

\end{document}
